\section{Conclusion}

The estimated coefficient revealed statistically significant and large conditional reading score gaps for learner that experienced bullying Monthly and Weekly, with more severe gaps among learners with Weekly bullying. South African coefficients experienced attenuation as the learner, school and household characteristics were controlled reflecting the observable differences between learners. The Brazilian coefficients remained stable suggesting that the included covariates account for less of the association in the sample. The interaction models revealed that, although gender may influence exposure to bullying, it did not significantly change the association between bullying frequency and reading scores. However, the country-specific home SES interactions showed contrasting patterns: higher SES intensified the monthly and weekly bullying penalties in South Africa but attenuated the weekly penalty in Brazil, while the Brazilian monthly interaction was not statistically significant. Nevertheless, the nature of cross section analysis does not address reverse causality, nor omitted variable bias. Subject to these limitations, the evidence suggests that bullying is associated with lower reading achievement and should be recognised as an important educational concern.

